% Unit 075 English translation of chapter6.tex lines 740--868.
% Source: Methods of Algebra, Volume 2, commit
% 9a5803ff77dd3257484cb177f851a73770a59dd3 (CC BY 4.0).
% stable-unit-id: o014.aljabr2.chapter6.induced-modules
% Coverage: Section 6.4 in full.

% segment-id: o014.li2.chapter6.induced-modules.h001
\section{Induced modules}\label{sec:induced-module}
\hypertarget{o014.aljabr2.chapter6.induced-modules}{ }

% segment-id: o014.li2.chapter6.induced-modules.p001
For a group homomorphism $\varphi:H\to G$,
Proposition~\ref{prop:pullback-two-adjoint} gives the right adjoint
$\Hom_H(\Bbbk[G],\cdot)$ and the left adjoint
$\Bbbk[G]\dotimes{\Bbbk[H]}(\cdot)$ of $\varphi^*$. This section focuses
on the common case in which $H\subset G$ is a subgroup and $\varphi$ is the
inclusion homomorphism. We begin by introducing notation.

% segment-id: o014.li2.chapter6.induced-modules.d001
\begin{definition}[Induced module]\label{def:induced-module}
	\index{$G$-module!induced}
	\index[sym1]{Ind-ind@$\Ind^G_H$, $\iInd^G_H$}
	Let $H$ be a subgroup of $G$. For an $H$-module $N$, define the
	$G$-modules
	\[ \Ind^G_H(N):=\Hom_H(\Bbbk[G],N),\qquad
	\iInd^G_H(N):=\Bbbk[G]\dotimes{\Bbbk[H]}N. \]
	As $N$ varies, both constructions give functors
	$H\dcate{Mod}\to G\dcate{Mod}$.
\end{definition}

% segment-id: o014.li2.chapter6.induced-modules.p002
Both operations may be called \emph{induction} from $H$ to $G$ in a broad
sense. The functor $\iInd^G_H$ is also sometimes called induction with finite
support; this terminology will be explained by
Lemma~\ref{prop:Ind-as-mappings} below.

% segment-id: o014.li2.chapter6.induced-modules.t001
\begin{proposition}\label{prop:Ind-preservation}
	The functors $\Res^G_H:G\dcate{Mod}\to H\dcate{Mod}$ and
	$\Ind^G_H,\iInd^G_H:H\dcate{Mod}\to G\dcate{Mod}$ are all exact.
	The functor $\Ind^G_H$ preserves injective objects,
	$\iInd^G_H$ preserves projective objects, and $\Res^G_H$ preserves both.
\end{proposition}
\begin{proof}
	We already know that $\Res^G_H$ is exact. Choosing coset representatives
	shows that $\Bbbk[G]$ is a free left $\Bbbk[H]$-module, so
	$\Hom_{\Bbbk[H]}(\Bbbk[G],\cdot)$ is exact. Similarly,
	$\Bbbk[G]$ is a free right $\Bbbk[H]$-module, so
	$\Bbbk[G]\dotimes{\Bbbk[H]}(\cdot)$ is exact. By the adjunctions, the
	assertions about preserving injective and projective objects follow
	directly from Proposition~\ref{prop:adjoint-injective-projective}.
\end{proof}

% segment-id: o014.li2.chapter6.induced-modules.p003
Proposition~\ref{prop:Ind-preservation} provides a way to construct injective
(or projective) resolutions in $G\dcate{Mod}$: take an injective (or
projective) resolution of the given $G$-module $M$ in
$\Bbbk\dcate{Mod}$, and then induce that resolution to $G$ using
$\Ind^G_{\{1\}}$ (or $\iInd^G_{\{1\}}$).

% segment-id: o014.li2.chapter6.induced-modules.p004
The following result is commonly called \emph{Shapiro's lemma}.

% segment-id: o014.li2.chapter6.induced-modules.t002
\begin{theorem}[B.\ Eckmann, D.\ K.\ Faddeev, A.\ Shapiro]
\label{prop:Shapiro}
	\index{Shapiro's lemma}
	Let $H$ be a subgroup of $G$. For every $H$-module $N$ and
	$n\in\Z$, there are canonical isomorphisms
	\[ \Hm^n\left(G,\Ind^G_H(N)\right)\simeq\Hm^n(H,N),\qquad
	\Hm_n\left(G,\iInd^G_H(N)\right)\simeq\Hm_n(H,N). \]
	The two sides of each isomorphism, viewed as functors in $N$, have long
	exact sequences, and the isomorphisms are compatible with those
	sequences.
\end{theorem}
\begin{proof}
	We give two methods. For cohomology, the first considers the functors
	\[ A^n:=\Hm^n\left(G,\Ind^G_H(\cdot)\right). \]
	Since $\Ind^G_H$ is exact, this family of functors has long exact
	sequences inherited from $\Hm^n(G,\cdot)$ and hence is a cohomological
	$\delta$-functor in the sense of Definition~\ref{def:delta-functor}.
	For $n=0$, Corollary~\ref{prop:Ind-invariant} gives a canonical
	isomorphism
	\[ A^0(N)=(\Ind^G_H(N))^G\rightiso N^H,\qquad
	\varphi\mapsto\varphi(1_G). \]
	To extend this canonically to an isomorphism of cohomological
	$\delta$-functors $A^n\rightiso\Hm^n(H,\cdot)$, it suffices to prove
	that $A^n$ is effaceable for $n>0$
	(Definition~\ref{def:erasable}, Proposition~\ref{prop:erasable-univ-delta}).
	By Proposition~\ref{prop:Ind-preservation}, if $N$ is embedded into an
	injective $H$-module $I$, then
	$\Ind^G_H(N)\hookrightarrow\Ind^G_H(I)$ and $\Ind^G_H(I)$ is an
	injective $G$-module. Hence $n>0$ implies
	$\Hm^n(G,\Ind^G_H(I))=0$, proving effaceability.

	The idea for homology is similar and uses the coeffaceability of
	homological $\delta$-functors. The required degree-zero isomorphism
	$N_H\rightiso(\iInd^G_H(N))_G$ also comes from
	Corollary~\ref{prop:Ind-invariant}.

	The second method works in the derived category. For cohomology, we again
	start from the canonical isomorphism
	\[ \underbracket{\Hom_G(\Bbbk,\Ind^G_H(\cdot))}_{
	\simeq\Ind^G_H(\cdot)^G}
	\simeq
	\underbracket{\Hom_H(\Bbbk,\cdot)}_{\simeq(\cdot)^H}
	:H\dcate{Mod}\to\Bbbk\dcate{Mod}. \]
	Both sides are left-exact functors. Since $\Ind^G_H$ preserves injective
	objects, taking right-derived functors in
	Theorem~\ref{prop:localization-triangulated-composite}~(iii) gives
	\begin{equation*}
		\RHom_G(\Bbbk,\mathrm{R}\Ind^G_H(\cdot))
		\simeq\RHom_H(\Bbbk,\cdot).
	\end{equation*}

	But $\Ind^G_H$ is exact, so $\mathrm{R}\Ind^G_H(\cdot)$ is determined by
	the universal property of localization and may still be denoted by
	$\Ind^G_H$. Thus at the level of derived categories there is an
	isomorphism of triangulated functors
	\[ \RHom_G(\Bbbk,\Ind^G_H(\cdot))
	\simeq\RHom_H(\Bbbk,\cdot). \]
	Taking $\Hm^n$ on both sides gives the desired result.

	The idea for homology is similar; the key is the isomorphism
	\[ \Bbbk\otimesL[{\Bbbk[G]}]
	\left(\Bbbk[G]\otimesL[{\Bbbk[H]}](\cdot)\right)
	\simeq\Bbbk\otimesL[{\Bbbk[H]}](\cdot). \]
	This is also a special case of Example~\ref{eg:change-of-rings} or of
	Proposition~\ref{prop:otimesL-associativity-constraint} on the
	associativity constraint. Since $\iInd^G_H$ is exact,
	$\Bbbk[G]\otimesL[{\Bbbk[H]}](\cdot)$ may also be replaced by
	$\Bbbk[G]\dotimes{\Bbbk[H]}(\cdot)$; the remaining steps are the same.
\end{proof}

% segment-id: o014.li2.chapter6.induced-modules.p005
Notice that if $\Res^G_H$ is interpreted as
$\Bbbk[G]\dotimes{\Bbbk[G]}(\cdot)$ according to
\eqref{eqn:G-module-pullback-tensor}, then, since
$\Res^G_H\Bbbk=\Bbbk$, the isomorphism
$\RHom_G(\Bbbk,\mathrm{R}\Ind^G_H(\cdot))
\simeq\RHom_H(\Bbbk,\cdot)$ reduces to the adjunction between $\RHom$ and
$\otimesL$ (Theorem~\ref{prop:RHom-otimesL-adjunction}).

% segment-id: o014.li2.chapter6.induced-modules.p006
Group cohomology and homology can be described concretely through the standard
complex and the standard chain complex of
Propositions~\ref{prop:groupcoh-cochain} and \ref{prop:grouph-chain}. We want
to make the isomorphisms in Theorem~\ref{prop:Shapiro} explicit in terms of
these complexes. There is no essential difficulty. To distinguish the group,
for every group $H$ write the complex of $H$-modules $\mathsf{L}$ from
Definition~\ref{def:resolution-trivial-module} as $\mathsf{L}^H$; it gives a
free resolution $\mathsf{L}^H\to\Bbbk$.

% segment-id: o014.li2.chapter6.induced-modules.l001
\begin{lemma}\label{prop:Shapiro-alpha-beta}
	Let $\varphi:H\to G$ be a group homomorphism. The corresponding
	homomorphism $H^{n+1}\to G^{n+1}$ induces a quasi-isomorphism of complexes
	$\mathsf{L}^H\to\varphi^*(\mathsf{L}^G)$ whose composite with
	$\varphi^*(\mathsf{L}^G)\to\varphi^*\Bbbk=\Bbbk$ equals
	$\mathsf{L}^H\to\Bbbk$.

	In particular, if $H$ is a subgroup of $G$, there is a quasi-isomorphism
	$\mathsf{L}^H\to\Res^G_H\mathsf{L}^G$ with the above composite property.
\end{lemma}
\begin{proof}
	It follows directly from the definitions that
	$\mathsf{L}^H\to\varphi^*(\mathsf{L}^G)$ is a morphism of complexes; the
	assertion about its composite is equally clear. Since $\varphi^*$ is
	exact, while $\mathsf{L}^G\to\Bbbk$ and $\mathsf{L}^H\to\Bbbk$ are both
	quasi-isomorphisms, $\mathsf{L}^H\to\varphi^*(\mathsf{L}^G)$ is a
	quasi-isomorphism as well.
\end{proof}

% segment-id: o014.li2.chapter6.induced-modules.t003
\begin{proposition}\label{prop:Shapiro-std-coh}
	Let $H$ be a subgroup of $G$ and $N$ an $H$-module. Identify
	$\Hm^n(G,\Ind^G_H(N))$ (or $\Hm^n(H,N)$) with the $\Hm^n$ of the
	standard complex $C(G,\Ind^G_H(N))$ (or $C(H,N)$), as in
	Proposition~\ref{prop:groupcoh-cochain}. Then the isomorphism
	\[ \Hm^n\left(G,\Ind^G_H(N)\right)\simeq\Hm^n(H,N),
	\qquad n\in\Z_{\geq0}, \]
	is induced by the morphism of complexes
	\begin{align*}
		C^m(G,\Ind^G_H(N))&\to C^m(H,N),\\
		[f:G^m\to\Ind^G_H(N)]&\mapsto
		[(h_1,\ldots,h_m)\mapsto f(h_1,\ldots,h_m)(1_G)]\\
		&\quad=(f|_{H^m})(\cdot)(1_G),
	\end{align*}
	where $h_1,\ldots,h_m\in H$.
\end{proposition}
\begin{proof}
	Again there are two methods. First, verify that the displayed maps
	$C^m(G,\Ind^G_H(N))\to C^m(H,N)$ do indeed form a morphism of complexes.
	Because $\Ind^G_H$ is exact, as $N$ varies these maps induce a morphism
	between cohomological $\delta$-functors. In degree zero, this map is
	clearly the canonical isomorphism
	$\Ind^G_H(N)^G\rightiso N^H$, $\varphi\mapsto\varphi(1_G)$. We already
	know that $\Hm^n(G,\Ind^G_H(\cdot))$ is a universal cohomological
	$\delta$-functor (Theorem~\ref{prop:Shapiro}). Thus the above morphism is
	unique and must agree with the isomorphism
	$\Hm^n(G,\Ind^G_H(N))\simeq\Hm^n(H,N)$ in that theorem.

	For the second method, let $M$ (respectively, $N$) be a bounded-below
	complex of $G$-modules (respectively, $H$-modules). Extend the functors
	$\Ind^G_H$ and $\Res^G_H$ degreewise to complexes, keeping the same
	notation. The adjunction of Proposition~\ref{prop:pullback-two-adjoint}
	lifts to $\Hom$ complexes:
	% Editorial correction: the Chinese source interchanges the H and G
	% subscripts on the two Hom complexes; the Indonesian witness restores
	% the adjunction with its correctly typed categories.
	\[ \Hom^\bullet_H(\Res^G_HM,N)
	\simeq\Hom^\bullet_G(M,\Ind^G_H(N)). \]
	A direct verification is not difficult; this can also be understood as
	the adjunction between $\Hom^\bullet$ and $\otimes$ described earlier
	(Proposition~\ref{prop:tensor-Hom-cplx}).

	Now let $N$ be an $H$-module and take an injective resolution
	$\iota:N\to I$. Then $\Ind^G_H(N)\to\Ind^G_H(I)$ is still an injective
	resolution, and naturality of the $\Hom$ complexes gives the following
	commutative diagram of complexes (the solid part):
	\begin{equation}\label{eqn:Shapiro-explicit-diagram}
	\begin{tikzcd}[row sep=small]
		\Hom^\bullet_H(\Bbbk,I) \arrow[d,"\alpha"'] &
		\Hom^\bullet_G(\Bbbk,\Ind^G_H(I)) \arrow[l,"\sim"'] \arrow[d] \\
		\Hom^\bullet_H(\Res^G_H\mathsf{L}^G,I) \arrow[d,"\beta"'] &
		\Hom^\bullet_G(\mathsf{L}^G,\Ind^G_H(I)) \arrow[l,"\sim"] \\
		\Hom^\bullet_H(\mathsf{L}^H,I) & \\
		\Hom^\bullet_H(\mathsf{L}^H,N) \arrow[u] &
		\Hom^\bullet_G(\mathsf{L}^G,\Ind^G_H(N))
		\arrow[uu] \arrow[dashed,l,"\gamma"]
	\end{tikzcd}
	\end{equation}
	where $\alpha$ and $\beta$ come from the morphism in
	Lemma~\ref{prop:Shapiro-alpha-beta}. Thus $\beta\alpha$ is the familiar
	map $\Hom^\bullet_H(\Bbbk,I)\to\Hom^\bullet_H(\mathsf{L}^H,I)$. The
	solid horizontal isomorphisms come from the adjunction, and all vertical
	arrows are quasi-isomorphisms.

	Recall the second proof of Theorem~\ref{prop:Shapiro}. The first row of
	\eqref{eqn:Shapiro-explicit-diagram} gives the derived-category form of
	\[ \Hm^n(H,N)\simeq\Hm^n(G,\Ind^G_H(N)),\qquad n\in\Z. \]
	On the other hand, the substance of
	Proposition~\ref{prop:groupcoh-cochain} is the isomorphisms of complexes
	\[ \Hom^\bullet_H(\mathsf{L}^H,N)\simeq C(H,N),\qquad
	\Hom^\bullet_G(\mathsf{L}^G,\Ind^G_H(N))
	\simeq C(G,\Ind^G_H(N)). \]
	It therefore remains only to supply the arrow $\gamma$ in
	\eqref{eqn:Shapiro-explicit-diagram} so that the square commutes, and to
	describe it. Since $\iota$ is monic, such a $\gamma$, if it exists, is
	unique.

	Take $m\in\Z_{\geq0}$ and
	$\phi\in\Hom^m_G(\mathsf{L}^G,\Ind^G_H(N))$. Inspecting the construction
	shows that the image of $\phi$ in $\Hom^m_H(\mathsf{L}^H,I)$ is
	\[ (h_0,\ldots,h_m)\mapsto
	\iota\left(\phi(h_0,\ldots,h_m)(1_G)\right). \]
	Comparison with the isomorphism in the proof of
	Proposition~\ref{prop:groupcoh-cochain} immediately shows that $\gamma$
	can be defined as asserted.
\end{proof}

% segment-id: o014.li2.chapter6.induced-modules.p007
We now state the homological version, using the notation for the standard
chain complex from Proposition~\ref{prop:grouph-chain}.

% segment-id: o014.li2.chapter6.induced-modules.t004
\begin{proposition}\label{prop:Shapiro-std-ho}
	Let $H$ be a subgroup of $G$ and $N$ an $H$-module. Identify
	$\Hm_n(G,\iInd^G_H(N))$ (or $\Hm_n(H,N)$) with the $\Hm_n$ of the
	standard chain complex $C(G,\iInd^G_H(N))$ (or $C(H,N)$), as in
	Proposition~\ref{prop:grouph-chain}. Then the isomorphism
	% Editorial correction: the Chinese source has Ind rather than iInd in
	% this displayed homology isomorphism; the surrounding statement,
	% theorem, and Indonesian witness all require iInd.
	\[ \Hm_n\left(G,\iInd^G_H(N)\right)\simeq\Hm_n(H,N),
	\qquad n\in\Z_{\geq0}, \]
	is induced by the morphism of complexes
	\begin{align*}
		C_m(H,N)&\to C_m(G,\iInd^G_H(N)),\\
		(h_1|\cdots|h_m)^\dagger\otimes y
		&\mapsto(h_1|\cdots|h_m)^\dagger\otimes(1_G\otimes y),
	\end{align*}
	where $h_1,\ldots,h_m\in H$ and $y\in N$.
\end{proposition}
\begin{proof}
	We explain only the derived-category approach. Take a projective
	resolution $\pi:P\to N$ in $H$-modules. The commutative diagram
	\eqref{eqn:Shapiro-explicit-diagram} from the proof of
	Proposition~\ref{prop:Shapiro-std-coh} must be replaced by
	\[\begin{tikzcd}[row sep=small]
		\Bbbk\dotimes{\Bbbk[H]}P \arrow[r,"\sim"] &
		\Bbbk\dotimes{\Bbbk[G]}\iInd^G_H(P) \\
		\Res^G_H\mathsf{L}^G\dotimes{\Bbbk[H]}P
		\arrow[r,"\sim"] \arrow[u,"\alpha"] &
		\mathsf{L}^G\dotimes{\Bbbk[G]}\iInd^G_H(P)
		\arrow[u] \arrow[dd] \\
		\mathsf{L}^H\dotimes{\Bbbk[H]}P \arrow[u,"\beta"] \arrow[d] & \\
		\mathsf{L}^H\dotimes{\Bbbk[H]}N \arrow[dashed,r,"\gamma"'] &
		\mathsf{L}^G\dotimes{\Bbbk[G]}\iInd^G_H(N),
	\end{tikzcd}\]
	where the solid horizontal arrows come from the associativity constraint
	for tensor products at the level of complexes, while all vertical arrows
	are quasi-isomorphisms. Once again, it remains to describe $\gamma$. The
	rest of the argument is similar to the cohomological version.
\end{proof}
